% Nouveau EGA V (cote 25): archive cote 25, site page 168 (PDF page 169).
% typescript. Draft compared: none.
% AI-transcribed by Claude Opus 5.5 (claude-opus-5-5, adaptive thinking, effort high), Message Batches API, batch ega25b (msgbatch_01A7n5YuXubUPNK1zsepi8RG), 2026-11-13; not yet checked by eye.
\scanpage{25/168}
\mspage{\hand{3}}

EGA IV ..... Ensemble singulier et\var{\xout{x}}{}{typewriter strike-out}{ega25_25_168-v0001} supersingulier.

Il me semble que le changement de point de vue possible signalé dans ma
dernière lettre peut-être récupéré après-coup, et qu'il peut y avoir intérêt
à partir d'abord de la situation où on a  un homomorphisme $\underline{E}\hand{\to}\underline{O}_X$ , d'où
un $V^o = V(u) = V(u(\underline{E})\underline{O}_X)$, etc. Une fois cette étude faite, on fera les
remarques suivantes: supposons donné un X lisse(de dimension relative \hand{d} ..)
sur Y, un sous-préschéma $V^o = V(\underline{J})$ de X, tel que $\underline{J}$\shortstack[l]{\var{\xout{$\hand{/\underline{J}^2}\unsure{=}N\unsure{\hand{\otimes}}$}}{}{typewriter strike-out}{ega25_25_168-v0002}\\\var{\xout{$K\hand{\hand{\otimes}}\unsure{O}$}}{}{typewriter strike-out}{ega25_25_168-v0003}} soit engendrab-%
le localement \var{\inserted{en les \var{\deleted{\illegible{1}}}{}{deletion}{ega25_25_168-v0005} $x\in V^o$}}{en les \var{\deleted{\illegible{1}}}{}{deletion}{ega25_25_168-v0005} $x\in V^o$}{insertion}{ega25_25_168-v0004} par m éléments ($m\hand{\geqslant}\unsure{0|1}$ étant donné), où ce qui revient au
même, tel que $V^o$ soit localement de présentation finie sur X (ou encore sur
Y) et que \var{\inserted[typed]{$\underline{N} =$}}{$\underline{N} =$}{insertion}{ega25_25_168-v0006} $\underline{N}_{V^o/X} = \underline{J}/\underline{J}^2$  soit localement engendrable par m éléments.  L'hy-%
pothèse peut aussi s'exprimer en disant que localement sur X, $\underline{J}$ peut être
représenté comme image d'un homomorphisme  u: $\underline{E}\hand{\to}\underline{O}_X$, $\underline{E}$ étant localement
libre (ou même libre, tant qu'on \unsure{r|m}este dans le local !) de rang m.  En d'aut\unsure{x|r}
res termes, localement on est dans la situation de départ envisagée dans
les pages 14 et suivantes. Utilisant une telle représentation, on construit
alors des \var{\inserted{sous-}}{sous-}{insertion}{ega25_25_168-v0007}préschémas \var{\inserted{fermés}}{fermés}{insertion}{ega25_25_168-v0008} $V^1 \hand{\subset} V^o$ , $V'' \hand{\subset} V^1$, $V^2 \hand{\subset} V^{\hand{1}}-V''$ . Il faut alors noter
(en proposition en forme) que ces sous-préschémas (et à fortiori les ensem-%
bles sous-jacents) sont en fait indépendants du choix de u:$\underline{E}\hand{\to}\underline{O}_X$ définiss-%
ant $V^o$, \var{\deleted{mais}}{}{deletion}{ega25_25_168-v0009} \var{\inserted{et}}{et}{insertion}{ega25_25_168-v0010} dépendent seu\var{\xout{\unsure{r}}}{}{typewriter strike-out}{ega25_25_168-v0011}\var{\over{l}{l}}{l}{writing over}{ega25_25_168-v0012}ement de \var{\deleted{\var{\xout{X/Y}}{}{typewriter strike-out}{ega25_25_168-v0014}}}{}{deletion}{ega25_25_168-v0013} $V^o$ ( et bien sur de X/Y et de
m !).  En effet, considérons d'abord  l'homomorphisme
\[ \hand{(*^o)} \qquad \underline{N}\hand{\otimes}\underline{O}_{V^o}\hand{\to}\hand{\underline{\Omega}}^1_{X/Y}\hand{\otimes}\underline{O}_{V^o} \ , \]
d'où  un homomorphisme
\[ \hand{(*_m)} \qquad \hand{\bigwedge}^m\underline{N}\hand{\otimes}\underline{O}_{V^o}\hand{\to}\hand{\underline{\Omega}}^{\var{\over{\unsure{1}}{m}}{m}{writing over}{ega25_25_168-v0015}}_{X/Y}\hand{\otimes}\underline{O}_{V^o}\ ; \]
on constate aussitôt que $V^1$ est le préschéma des zéros de cet homomorphisme
(qui ``existe'' grâce au fait que $\hand{\underline{\Omega}}^{\var{\over{\unsure{1|m}}{m}}{m}{writing over}{ega25_25_168-v0016}}_{X/Y}\hand{\otimes}\underline{O}_{V^{\var{\over{\unsure{1}}{o}}{o}{writing over}{ega25_25_168-v0017}}}$  sur  $V^o$ est localement libre
de rang fini, et se définit comme réprésentant le foncteur qu'on devine
(Sch)$^o_{\hand{/}\var{\over{\unsure{V}}{V}}{V}{writing over}{ega25_25_168-v0018}^{\hand{o}}}$ $\hand{\longrightarrow}$ (Ens) ; c'est aussi le sous-préschéma de $\var{\over{V}{V}}{V}{writing over}{ega25_25_168-v0019}^{\hand{o}}$

\mspage{\hand{168}}
